\chapter{A\_json\_commander.py}

This is a script provided by Fraunhofer IPM, which can be used to send commands using the JSON interface to programs such as Pulsalyzer. This script thus enables a user to control programs (such as Pulsalyzer) programmatically.

\section{How It Works}
The script connects to a local TCP socket on 127.0.0.1:44444 (JSON\_Router) and routes JSON-formatted commands to programs and their endpoints.

For more details regarding the JSON interface, its features and detailed usage, visit the local URL \url{http://localhost:5000/ipm/help/jsoniface_help} after launching \textbf{JSONServer}, which is provided by Fraunhofer IPM.

\section{Key Function}
\texttt{send\_commands(cmds, routing, host="127.0.0.1", port=44444)}

\begin{table}[ht!]
    \centering
    \begin{tabular}{lll}
        \toprule
        Parameter & Type & Description                                                                  \\
        \midrule
        cmds      & list & List of command dictionaries (e.g., \{"cmd": "reprocess\_pointcloud", ...\}) \\
        routing   & list & Routing path to endpoint (see \autoref{ch:scriptF} for format)               \\
        host      & str  & IP address of the host (default: \texttt{"127.0.0.1"} (localhost))           \\
        port      & int  & TCP port of the JSON\_Router (default: 44444)                                \\
        \bottomrule
    \end{tabular}
\end{table}

\section{Usage}
This script is not run directly. It is to be imported by other scripts, thus controlling programs programmatically. [See \autoref{ch:scriptF} for an example]

\begin{lstlisting}[language=python, numbers=none]
import A_json_commander as json_commander
json_commander.send_commands(cmds, routing)
\end{lstlisting}

\textbf{Note:} \texttt{`A\_json\_commander.py'} must be placed in the same directory as any script that imports it.
